% Esquisse d'un programme (cote 119, copy 2): archive cote 119, site page 134 (PDF page 135), typescript page 56.
% typescript. Draft compared: esquisse/pages/esq_119_058.tex (copy 1, site 58).
% AI-transcribed by Claude Opus 5.5 (claude-opus-5-5, adaptive thinking, effort high), Message Batches API, batch pilotnext (msgbatch_01A9sdu8SYq23bEtqX9tUncQ), 2026-11-13; not yet checked by eye.
\scanpage{119/134}
\mspage{56}

fondements de topologie utilisés pour la formuler n'excluaient
pas les phénomènes de sauvagerie. Il va (je l'espère) sans dire
que la nécessité de développer de nouveaux fondements pour la
topologie ``géométrique'' n'exclut nullement que les phénomènes en
question, comme toute chose sous le ciel, ont leur raison d'être
et leur propre beauté. Des fondements plus adéquats ne suppri-%
meront pas ces phénomènes, mais nous permettront de les situer à
leur juste place, comme des ``cas limites'' de phénomènes de ``vraie''
topologie..

($^7$) En fait, pour reconstituer ce système d'espaces
\[ (i_o,\ldots,i_n) \longmapsto X_{i_o,\ldots,i_n} \]
contravariant sur Drap(I) (pour l'inclusion des drapeaux), il
suffit de connaître les $X_i$ (ou ``\uline{strates déployées}'') et les $X_{i,j}$
(ou ``\uline{tubes de raccord}'') pour $i,j \in I$, $i < j$, et les morphismes
$X_{i,j} \longrightarrow X_i$ (qui sont des inclusions ``bordantes'') et $X_{i,j} \rightarrow X_i$
(qui sont des fibrations propres, dont les fibres $F_{i,j}$ sont
appelées ``\uline{fibres de raccord}'' pour les strates d'indice $i$ et $j$).
Dans le cas d'une multiplicité modérée cependant, il faut connaî-%
tre de plus les ``\uline{espaces de jonction}'' $X_{i,j,k}$ $(i < j < k)$ et ses
morphismes dans $X_{i,j}$ , $X_{j,k}$ , et surtout $X_{i,k}$ , s'insérant dans
le diagramme commutatif hexagonal suivant, où les deux carrés de
droite sont cartésiens, les flèches $\hookrightarrow$ sont des immersions (pas
nécessairement des plongements ici), et les autres flèches sont
des fibrations propres :

\diagram{
\begin{tikzcd}
 & X_{ik} \arrow[rr, hook] \arrow[dl] & & X_k & \\
X_i & & X_{ijk} \arrow[ul, hook] \arrow[rr, hook] \arrow[dl] & & X_{jk} \arrow[ul, hook] \arrow[dl] \\
 & X_{ij} \arrow[ul] \arrow[rr, hook] & & X_j &
\end{tikzcd}
}

(NB. Ce diagramme définit $X_{ijk}$ en termes de $X_{ij}$ et $X_{jk}$ sur $X_j$,
mais non la flèche $X_{ijk} \longrightarrow X_{ik}$ , car $X_{ik} \longrightarrow X_k$ n'est \uline{pas} né-%
cessairement un plongement.)

\margin{134}
